\clearpage
\VerdierPhysicalPageAnchor{0037}%
\VerdierSemanticAnchor{ch1-1-2-1-coherence}
并且对 $G$ 中每一组三个元素 $(g,h,k)$，都有关系式

\VerdierSemanticAnchor{ch1-1-2-1-eq3}
\begin{equation}
  \bigl(T(g)\star c(k,h)\bigr)c(kh,g)
  =\bigl(c(h,g)\star T(k)\bigr)c(k,hg).
  \tag{1.2.1.3}
\end{equation}

\VerdierSemanticAnchor{ch1-1-2-2}
\noindent\textbf{定义 1.2.2.}\quad
设 $G$ 为一个群。一个 $G^\circ$-\emph{范畴}（其中 $G^\circ$ 是 $G$ 的反群）
是一个范畴 $\mathcal C_e$，连同一族函子 $T(g)$（$g\in G$）和一族函子同构
$c(h,g)$（$g,h\in G$），满足性质 (1.2.1.1)、(1.2.1.2) 与 (1.2.1.3)。

\VerdierSemanticAnchor{ch1-1-2-2-transition}
因此，一个 $G^\circ$-范畴就是一个范畴 $\mathcal C_e$，配备群 $G^\circ$ 的一个
同构意义下作用。我们以 $(\mathcal C_e,G^\circ,T,c)$ 表示这个
$G^\circ$-范畴；同构 $c(h,g)$ 称为\emph{过渡同构}。

\VerdierSemanticAnchor{ch1-1-2-2-strict}
当过渡同构都是恒等同构时——这特别蕴含对每一对 $g,h\in G$ 都有
$T(hg)=T(g)T(h)$——称这个 $G^\circ$-范畴是\emph{严格的}。因此，这样一个对象
由 $G^\circ$ 在范畴 $\mathcal C_e$ 的自同构群中的一个表示所确定。

\VerdierSemanticAnchor{ch1-1-2-3}
\noindent\textbf{1.2.3.}\quad
现在设 $\mathcal C$ 与 $\mathcal C'$ 为两个 $G$-分次范畴，并设
$F\colon\mathcal C\to\mathcal C'$ 为一个分次函子。在 $\mathcal C$ 与
$\mathcal C'$ 上选取正规劈分。这些劈分的选择确定函子同构

\VerdierSemanticAnchor{ch1-1-2-3-eq1}
\begin{equation}
  m(g)\colon F_eT(g)\longrightarrow T'(g)F_e,\qquad
  g\in G,\qquad m(e)=\text{$F_e$ 上的恒等同构}
  \tag{1.2.3.1}
\end{equation}
使得对每一对 $g,h\in G$，下图交换：

\VerdierSemanticAnchor{ch1-1-2-3-eq2}
\begin{equation}
\begin{tikzcd}[column sep=huge,row sep=large]
F_eT(hg)
  \arrow[r,"m(hg)"]
  \arrow[d,"{F_e\star c(h,g)}",swap]
&
T^{\prime}(hg)F_e
  \arrow[dd,"{c^{\prime}(h,g)\star F_e}"]
\\
F_eT(g)T(h)
  \arrow[d,"m(g)\star T(h)",swap]
&
{}
\\
T^{\prime}(g)F_eT(h)
  \arrow[r,"T^{\prime}(g)\star m(h)"]
&
T^{\prime}(g)T^{\prime}(h)F_e\mathpunct{.}
\end{tikzcd}
\tag{1.2.3.2}
\end{equation}

\VerdierSemanticAnchor{ch1-1-2-4}
\noindent\textbf{定义 1.2.4.}\quad
设 $(\mathcal C_e,G^\circ,T,c)$ 与
$(\mathcal C'_e,G^\circ,T',c')$ 为两个 $G^\circ$-范畴。一个
$G^\circ$-\emph{函子}是一个函子

\VerdierSemanticAnchor{ch1-1-2-4-fe}
\[
  F_e\colon\mathcal C_e\longrightarrow\mathcal C'_e
\]
